\section{Positive maps and channels}

\begin{definition}[Positive map]\label{def:positive_map}
    \lean{IsPositiveMap}
    \leanok
    A linear map $E : \MN{n} \to \MN{m}$ is \emph{positive}
    if $E(X) \ge 0$ whenever $X \ge 0$, where $X \ge 0$
    means that $X$ is positive semidefinite.
\end{definition}

\begin{remark}\leanok
    We use the Loewner order on matrices: $X \le Y$ means
    $Y - X \ge 0$.
\end{remark}
\begin{definition}[Completely positive map]\label{def:cp_map}
    \lean{IsCPMap}
    \leanok
    A linear map $E : \MN{D} \to \MN{D}$ is \emph{completely positive} (CP)
    if it admits a \emph{Kraus representation}: there exist operators
    $\{K_i\}_{i=0}^{r-1}$ with $K_i \in \MN{D}$ such that, for every
    $X \in \MN{D}$,
    \[
        E(X) = \sum_{i=0}^{r-1} K_i X K_i^\dagger.
    \]
    By Choi's theorem, this is equivalent
    to $E \otimes \id_n$ being positive for all~$n$.
\end{definition}

\section{Irreducibility}

\begin{definition}[Orthogonal projection]
    \label{def:orthogonal_projection_channel}
    \lean{IsOrthogonalProjection}
    \leanok
    A matrix $P \in \MN{D}$ is an \emph{orthogonal projection} if
    $P = P^\dagger$ and $P^2 = P$.
\end{definition}

\begin{definition}[Irreducible map]\label{def:irreducible_cp}
    \lean{IsIrreducibleMap}
    \leanok
    \uses{def:orthogonal_projection_channel}
    A linear map
    $E : \MN{D} \to \MN{D}$ is \emph{irreducible}
    if whenever $P$ is an orthogonal projection satisfying
    $E(P \MN{D} P) \subseteq P \MN{D} P$,
    then $P = 0$ or $P = \Id$.
    This is \cite[Theorem~6.2(1)]{Wolf2012Quantum}.
    The definition applies to any linear map; complete positivity
    is not required.
\end{definition}

\section{Peripheral spectrum and primitivity}

\begin{definition}[Peripheral spectrum]\label{def:peripheral_spectrum}
    \lean{peripheralSpectrum}
    \leanok
    The \emph{peripheral spectrum} of a bounded linear operator $T$
    is the set of spectral values $\mu$ with
    $|\mu|$ equal to the spectral radius of $T$.
\end{definition}

\begin{definition}[Peripheral eigenvalues]\label{def:peripheral_eigenvalues}
    \lean{peripheralEigenvalues}
    \leanok
    The \emph{peripheral eigenvalues} of a linear map $E$ on
    a finite-dimensional space
    are the eigenvalues $\lambda$ on the unit circle:
    $|\lambda| = 1$.
    For a channel, the spectral radius is~$1$
    \cite[Proposition~6.1]{Wolf2012Quantum}, so these coincide with
    the eigenvalues of maximal modulus.

    The condition $|\lambda| = 1$ is appropriate for channels, since the
    spectral radius of a channel is~$1$; for a general linear map with
    spectral radius $r \neq 1$ the condition would generalise to
    $|\lambda| = r$.
\end{definition}

\begin{remark}[Peripheral conventions for channels]%
    \label{rem:peripheral_conventions_channels}
    \leanok
    \uses{def:peripheral_spectrum, def:peripheral_eigenvalues}
    The spectral-radius convention of
    Definition~\ref{def:peripheral_spectrum} is the usual bounded-operator
    notion. For channels the spectral radius is $1$, so the channel arguments
    below use the unit-circle set of Definition~\ref{def:peripheral_eigenvalues}.
\end{remark}

\begin{definition}[Primitive map]\label{def:primitive_channel}
    \lean{IsPrimitive}
    \leanok
    \uses{def:peripheral_eigenvalues}
    A linear map is \emph{primitive} if its only peripheral eigenvalue
    is $\lambda = 1$, i.e.\
    $\mathrm{peripheral}(E) = \{1\}$.
    For channels this corresponds to \cite[Theorem~6.7]{Wolf2012Quantum}.
    This definition applies to any linear endomorphism, not only
    to channels.
\end{definition}
